\documentclass[11pt]{article}
\usepackage[OT1]{fontenc}
\renewcommand{\rmdefault}{cmr}
\renewcommand{\sfdefault}{cmss}
\renewcommand{\ttdefault}{cmtt}
\usepackage{amsmath,amssymb,amsthm,mathtools}
\usepackage[margin=27mm]{geometry}
\usepackage{microtype,booktabs,longtable,array}
\usepackage[hidelinks]{hyperref}
\setlength{\emergencystretch}{2em}
\setlength{\parskip}{3pt plus 1pt}
\setcounter{tocdepth}{1}
\renewcommand{\labelitemi}{$\bullet$}
\newtheorem{theorem}{Theorem}[section]
\newtheorem{lemma}[theorem]{Lemma}
\newtheorem{proposition}[theorem]{Proposition}
\newtheorem{corollary}[theorem]{Corollary}
\theoremstyle{definition}
\newtheorem{definition}[theorem]{Definition}
\newtheorem{example}[theorem]{Example}
\theoremstyle{remark}
\newtheorem{remark}[theorem]{Remark}
\newcommand{\N}{\mathbb N}
\newcommand{\Q}{\mathbb Q}
\newcommand{\R}{\mathbb R}
\newcommand{\Z}{\mathbb Z}
\newcommand{\supp}{\operatorname{supp}}
\newcommand{\Con}{\operatorname{Con}}
\newcommand{\ord}{\operatorname{ord}}
\newcommand{\dom}{\operatorname{dom}}
\newcommand{\sgn}{\operatorname{sgn}}
\newcommand{\res}{\operatorname{Res}}
\newcommand{\rad}{\operatorname{rad}}
\newcommand{\CK}{\mathsf{CK}}
\newcommand{\KT}{\mathsf{KT}}
\newcommand{\DCB}{\mathsf{DCB}}
\newcommand{\FS}{\mathsf{FS}}
\newcommand{\Ag}{\mathsf{Ag}}
\newcommand{\Know}{\mathsf K}
\newcommand{\Truth}{\mathsf T}
\newcommand{\Prv}{\mathsf{Pr}}
\newcommand{\code}[1]{\ulcorner #1\urcorner}
\newcommand{\Tot}{\operatorname{Tot}}
\numberwithin{equation}{section}

\title{Transformations Between Deterministic and Stochastic\\Reinforcement-Learning Frameworks}
\author{}
\date{28 October 2026\\Working manuscript}
\begin{document}
\maketitle
\begin{abstract}
Alexander and Pedersen define a transformation of reinforcement-learning
frameworks by preservation of pairwise agent rankings in each transformed
environment, together with two existential nontriviality conditions. We
classify transformations between their four standard deterministic and
stochastic frameworks under that definition. With at least two actions,
a transformation exists exactly when the source agents are deterministic
or the target agents are stochastic. The environment type does not affect
existence. Positive cases follow from a one-parameter policy family and
a three-valued choice of source environments. Negative cases follow from
a finite scripted environment separating deterministic policies and the
convexity of values attainable by stochastic policies. We also exhibit
finite abstract frameworks showing that the full transformation notion
need not compose, and distinguish it from a stronger cross-environment
condition. The maps in the classification may be noncomputable and
noninjective; the theorem concerns the stated mathematical definition,
not effective conversion of arbitrary learned policies.
\end{abstract}
\tableofcontents

\section{What should a change of reinforcement-learning framework preserve?}

A reinforcement-learning framework specifies agents, environments and
the value obtained when an agent interacts with an environment.
Different formalizations permit deterministic or randomized choices,
different reward sets, and different conventions for who acts first.
Whether such changes preserve an agent's relative ability is a natural
foundational question.

Alexander and Pedersen \cite{AP} propose an abstract definition based
on rankings. An agent map converts source agents into target agents;
an environment map converts target environments back into source
environments. Two converted agents must have the same strict ordering
in a target environment as their source versions have in its converted
source environment. This definition compares agents \emph{within one
environment}. It makes no demand about comparing one agent's scores
across different environments.

We give a complete classification for the four standard frameworks
in version~5 of their paper. A previous counterexample note
\cite{Counterexample} supplies one positive instance contradicting the
claimed deterministic-to-stochastic obstruction. Here a single positive
construction decides every positive case, and a separate obstruction
decides every remaining case. Both arguments are included in full.

\begin{theorem}[Classification]\label{thm:classification}
Assume the finite action alphabet has at least two elements. Under
Definition~\ref{def:transformation}, a transformation between the four
frameworks of Definition~\ref{def:four} exists if and only if
\[
 \text{the source agents are deterministic}
 \quad\text{or}\quad
 \text{the target agents are stochastic}.
\]
For a singleton action alphabet, no such transformation exists,
including from a framework to itself.
\end{theorem}

The singleton qualification comes from the definition's nontriviality
requirement, not from a technical failure of the proof. With only one
policy, no pair of image agents can have different values.

\section{Frameworks, interactions and the exact transformation condition}

An abstract framework is a triple $F=(A,E,V)$ with
$V:A\times E\to\R$. Agents and environments need not have effective
descriptions at this level.

\begin{definition}\label{def:transformation}
A transformation from $F=(A,E,V)$ to
$\overline F=(\overline A,\overline E,\overline V)$ is a pair of maps
$f:A\to\overline A$, $g:\overline E\to E$ such that
\begin{equation}\label{eq:faithful}
 \overline V(f(p),\mu)<\overline V(f(q),\mu)
 \quad\Longleftrightarrow\quad
 V(p,g(\mu))<V(q,g(\mu))
\end{equation}
for every $p,q,\mu$, and such that:
\begin{enumerate}
\item some two image agents have strictly different values in one
      target environment;
\item some image agent has strictly different values in two target
      environments.
\end{enumerate}
\end{definition}

Equation~\eqref{eq:faithful} also preserves ties: apply it in both
orders and use the trichotomy of real order. It does not require $f$
or $g$ to be injective, computable, continuous, or faithful to numerical
values themselves. In particular the two nontriviality clauses are
existential requirements on the image, not global information-preservation
conditions.

\subsection{The four standard frameworks}

Fix finite disjoint alphabets $\mathcal A$ of actions and $\mathcal P$
of percepts. Each is nonempty. A reward function
$r:\mathcal P\to\Z$ has both 0 and 1 in its range. Choose percepts
$z_0,z_1$ with those rewards. Interaction is percept-first:
\[
 z^{(1)},a^{(1)},z^{(2)},a^{(2)},\ldots.
\]
An agent history is a finite alternating sequence ending in a percept;
an environment history is empty or ends in an action. A deterministic
policy gives one action at every agent history. A stochastic policy
gives a rational probability distribution on actions at every such
history. Environments similarly give a percept or a rational distribution
on percepts. A deterministic object is identified with its point-mass
stochastic version when necessary.

Policies are defined on all histories of the appropriate form,
including histories which would have probability zero in a particular
interaction. They have no effectiveness restriction.

Let $V_n(p,\mu)$ be the expected sum of rewards in the first $n$
percepts of the alternating process. An environment is \emph{admissible}
when $V(p,\mu)=\lim_n V_n(p,\mu)$ exists as a finite real for every
policy of either kind. There is no common normalization such as
$|V|\le1$. This condition is stronger than existence of a value just
against the policies occurring in the image of a proposed transformation.

\begin{definition}\label{def:four}
For $X,Y\in\{D,S\}$, write $F_{XY}$ for the framework with agent
type $X$ and environment type $Y$, using all agents of that type and
all admissible environments of that type. The first subscript always
denotes agents. $D$ means deterministic and $S$ stochastic.
\end{definition}

These are the four frameworks in \cite[Definition~9]{AP}. The finite
play convention is the alternating percept-first process stated there;
values here are defined directly by that process.

\section{The positive cases: one scalar of agent behavior is enough}

Choose distinct actions $a_0,a_1$. For a source policy $p$, define
\[
 t(p)=p(a_0\mid\langle z_0\rangle).
\]
For a deterministic source, interpret this as 1 if its first response
to $z_0$ is $a_0$, and 0 otherwise. Thus $t$ is binary in that case;
for stochastic sources it is rational in $[0,1]$. Both endpoints are
attained by constant-action policies.

For each rational $t\in[0,1]$, define a target policy $Q_t$ which, at
its first action, chooses $a_1$ with probability $t$ and $a_0$ with
probability $1-t$, independently of the initial percept. At every later
history it chooses $a_0$. This defines the policy on every history;
no persistent hidden random seed is required. The endpoints $Q_0,Q_1$
are deterministic.

\begin{lemma}\label{lem:affine}
For every admissible target environment $\mu$,
\[
 V(Q_t,\mu)=A_\mu+t(B_\mu-A_\mu),
 \qquad A_\mu=V(Q_0,\mu),\ B_\mu=V(Q_1,\mu).
\]
\end{lemma}
\begin{proof}
For each finite horizon, the trace law of $Q_t$ is the mixture of the
trace laws of $Q_0$ and $Q_1$ with weights $1-t,t$. The only varying
policy choice is the first action; after it all three policies have
the same prescribed continuation. Conditioning on the first percept
does not alter those weights. The environment may depend on the whole
history, but for a fixed trace its own probability factors are the
same in all three comparisons. Finite summation gives the identity for
$V_n$. Take the limit of this two-term linear combination. Both
endpoint limits are finite by admissibility.
\end{proof}

Three simple source environments suffice to match the possible
orderings in this affine family. Each emits $z_0$ initially and zero
rewards forever after its second percept:
\begin{itemize}
\item $e_+$ gives reward 1 at the second percept exactly when the
      first action was $a_0$;
\item $e_-$ gives reward 1 exactly when that first action was not $a_0$;
\item $e_0$ always gives reward zero.
\end{itemize}
Choose $z_1$ or $z_0$ to implement the reward. These rules specify the
environments on off-path histories as well. Their values against any
source policy are $t(p),1-t(p),0$, respectively. They are deterministic
and admissible against every deterministic or stochastic policy.

\begin{proposition}\label{prop:positive}
If source agents are deterministic or target agents are stochastic,
there is a transformation.
\end{proposition}
\begin{proof}
Set $f(p)=Q_{t(p)}$. Its codomain has the required type: either only
the deterministic endpoints are needed or all rational $Q_t$ are
allowed. For any target environment put
\[
 g(\mu)=
 \begin{cases}
 e_+,&B_\mu>A_\mu,\\
 e_-,&B_\mu<A_\mu,\\
 e_0,&B_\mu=A_\mu.
 \end{cases}
\]
Use the point-mass version if the source environment class is stochastic.
The affine identity shows that every image pair is ordered by increasing
$t$, decreasing $t$, or a constant value according to the same three
cases. The source environments realize exactly these orderings, so
\eqref{eq:faithful} holds for every pair, including ties, in every target
environment.

For the first nontriviality condition, take a target environment which
emits $z_0$, then rewards the first action $a_1$, then gives zeros.
It assigns values 0 and 1 to $Q_0,Q_1$, both of which occur in the
image. For the second, compare the always-zero target environment
with one which emits $z_0$, then $z_1$ independently of action, then
zeros. Every image policy has values 0 and 1 in these two environments.
All these environments are admissible in either environment type.
\end{proof}

The sign choice in $g$ is a set-theoretic comparison of two real values.
It need not be computable from any description of $\mu$. This is
allowed by the definition, and is central to the scope of the theorem.

\section{A finite environment which separates deterministic policies}

The negative argument needs more than two policies being syntactically
different. They can differ only at histories inconsistent with their
own previous actions and still have identical value in every environment.
An actual value distinction avoids that problem.

\begin{lemma}[Two-value separator]\label{lem:separator}
Suppose deterministic policies $\alpha,\beta$ have different values
in some admissible environment. There is an admissible deterministic
environment $\eta$ and an integer $C$ such that every deterministic
policy has value either $C$ or $C+1$ in $\eta$, with $\alpha$ and
$\beta$ taking different values.
\end{lemma}
\begin{proof}
There must be a finite prescribed percept sequence on which their
recursively generated actions first differ. Otherwise their actions
would agree along every prescribed percept sequence; induction on
trace length would then give identical finite trace distributions in
every environment, and consequently identical limiting values.

Choose a sequence $y_1,\ldots,y_k$ with first action divergence at
stage $k$. Let $b_1,\ldots,b_{k-1}$ be the common earlier actions and
$b_k$ the action taken by $\alpha$ at stage $k$.
Define $\eta$ to emit this entire percept script, independently of
actions. At percept $k+1$ emit $z_1$ if and only if the recorded action
sequence is exactly $b_1,\ldots,b_k$; otherwise emit $z_0$.
Emit $z_0$ thereafter.

Every deterministic policy receives the same script reward
$C=\sum_{i=1}^k r(y_i)$ and then either 0 or 1. In particular $\alpha$
receives $C+1$ and $\beta$ receives $C$. Every stochastic policy has
value between these endpoints. All nonzero rewards occur at a fixed
finite horizon, so admissibility holds against the full policy class.
\end{proof}

It matters that the script continues after an unexpected action.
Aborting it on deviation could change the script reward for third
policies and destroy the two-value conclusion. The offset $C$ also
allows prescribed percepts to carry arbitrary integer rewards; no
zero-reward copy of every observation is assumed.

\section{Mixing policies without a hidden random seed}

For stochastic source policies we need a policy whose value is the
midpoint of two given values in \emph{every} admissible environment.
Pointwise averaging of action distributions is generally insufficient:
it can redraw the branch at each step instead of mixing entire plays.
The correct behavioral policy updates the branch weights from the
recorded action history.

For an agent history $h$, let $L_p(h)$ be the product of the probabilities
that $p$ assigns to the earlier actions recorded in $h$, each evaluated
at its recorded preceding history. Do not multiply probabilities of
percepts. Let $0<w<1$ be rational. Define
\begin{equation}\label{eq:mixture}
 m(a\mid h)=
 \frac{wL_p(h)p(a\mid h)+(1-w)L_q(h)q(a\mid h)}
      {wL_p(h)+(1-w)L_q(h)}
\end{equation}
when the denominator is positive; otherwise use the point mass at $a_0$.
This is a rational distribution at every history.

\begin{lemma}\label{lem:mixture}
The policy $m$ satisfies
\[
 V(m,\nu)=wV(p,\nu)+(1-w)V(q,\nu)
\]
for every admissible environment $\nu$.
\end{lemma}
\begin{proof}
Induction over recorded actions gives
$L_m(h)=wL_p(h)+(1-w)L_q(h)$. For a positive denominator, multiplication
by \eqref{eq:mixture} cancels that denominator and gives the identity
at the next action. For a zero denominator both nonnegative summands
vanish, as do all their action extensions; the arbitrary off-path
choice cannot revive them. Extending a history by a percept leaves
the action-only products unchanged.

For each finite full trace in a fixed environment, its percept-likelihood
product is a common factor multiplying these three action products.
It may depend on the recorded actions; that does not matter because
the trace is fixed in this comparison. Finite trace probabilities,
and then finite reward expectations, satisfy the same mixture identity.
Take limits of the two-term linear combination. Admissibility supplies
the required finite limits for all three policies.
\end{proof}

\section{The negative cases and the complete table}

\begin{proposition}\label{prop:negative}
There is no transformation with stochastic source agents and
deterministic target agents, for either environment type.
\end{proposition}
\begin{proof}
Suppose $(f,g)$ were one. The first nontriviality condition supplies
source policies $p,q$ whose deterministic images have different values
in some target environment. Lemma~\ref{lem:separator} supplies a
deterministic environment $\eta$ giving every deterministic target
policy one of two values $C,C+1$, with $f(p),f(q)$ taking opposite
values. Its point-mass version is also available in a stochastic
target environment class.

Faithfulness makes $V(p,g(\eta))$ and $V(q,g(\eta))$ distinct.
Let $m$ be their midpoint policy from Lemma~\ref{lem:mixture}, with
$w=1/2$. Its source value is strictly between them. Applying
\eqref{eq:faithful} to the two comparisons forces $V(f(m),\eta)$
strictly between $C$ and $C+1$. No deterministic target policy has
such a value in $\eta$. This is a contradiction.
\end{proof}

The two propositions prove Theorem~\ref{thm:classification} when there
are at least two actions. If there is one action, both policy classes
have only one member, so the first nontriviality condition fails for
every map. The full table is therefore, for at least two actions,
\[
\begin{array}{c|cccc}
 \text{source}\backslash\text{target}&DD&DS&SD&SS\\\hline
 DD&\text{yes}&\text{yes}&\text{yes}&\text{yes}\\
 DS&\text{yes}&\text{yes}&\text{yes}&\text{yes}\\
 SD&\text{no}&\text{no}&\text{yes}&\text{yes}\\
 SS&\text{no}&\text{no}&\text{yes}&\text{yes}
\end{array}
\]
The environment type disappears from the criterion. Under this
definition there are two mutual-transformability classes, separated
by agent randomization, and a one-way transformation from the
deterministic class to the stochastic class. This concerns the
specific four frameworks; it is not a general composition theorem.

\section{Composition and a stronger comparison condition}

\subsection{The existential nontriviality conditions need not survive}

Faithfulness alone composes: substitute one pair of maps into the
other in \eqref{eq:faithful}. The full definition has an additional
issue: witnesses to nontriviality can fall outside the image of the
composite agent map or outside the selected environments.

\begin{proposition}\label{prop:noncomposition}
There are transformations $F\to G$ and $G\to H$ whose naturally
composed maps fail both nontriviality conditions.
\end{proposition}
\begin{proof}
The following value matrices give an example,
with agents as rows and environments as columns:
\[
 V=\begin{pmatrix}0&0\\1&0\end{pmatrix},\quad
 W=\begin{pmatrix}0&0&0\\1&0&0\\2&1&2\end{pmatrix},\quad
 Z=\begin{pmatrix}0&0\\0&0\\1&2\end{pmatrix}.
\]
From $V$ to $W$, map the two agents to rows 1 and 2, and map the
three target columns to source columns 1,2,2. On the selected rows
the entries agree exactly, so faithfulness holds. Column 1 separates
the image rows; image row 2 varies between columns 1 and 2.
From $W$ to $Z$, use the identity row map and select source columns
2,3. Again entries agree exactly. Row 3 witnesses both nontriviality
conditions. The composite selects the first two rows of $Z$, all
of whose entries are zero, so neither condition holds.

\end{proof}

This example addresses composition of the given maps. It does not
exclude the existence of a different transformation from $F$ to $H$.

\subsection{Comparing scores across environments changes the problem}

A stronger requirement would ask, for all $p,q,\mu,\nu$,
\begin{equation}\label{eq:cross}
 \overline V(f(p),\mu)<\overline V(f(q),\nu)
 \quad\Longleftrightarrow\quad
 V(p,g(\mu))<V(q,g(\nu)).
\end{equation}
Definition~\ref{def:transformation} is the special case $\mu=\nu$
and does not imply this condition. The difference has mathematical
consequences.

\begin{proposition}\label{prop:cross}
Assume source agents and environments are deterministic, with rewards
in a fixed lattice $d^{-1}\Z$, $d\ge1$, and convergent total rewards.
Assume the target has rational initial mixtures of any two admissible
environments. There is no pair satisfying \eqref{eq:cross} and the
second nontriviality condition.
\end{proposition}
\begin{proof}
A convergent sequence of source partial reward sums in $d^{-1}\Z$
is eventually constant. Every source value therefore lies in that
lattice. By \eqref{eq:cross}, all target values attained by image
agents have an order-preserving image in this lattice: equality is
also preserved by applying strict comparison in both orders.

The second nontriviality condition gives one image agent with values
$a<b$ in two target environments. Rational initial mixtures produce
values $a+t(b-a)$ for all rational $0<t<1$. These would map to
arbitrarily many distinct lattice points between the fixed lattice
images of $a$ and $b$, an impossibility.
\end{proof}

Initial mixtures can be realized behaviorally by the environment
analogue of \eqref{eq:mixture}, using past percept likelihoods.
For each policy, finite expectations are mixed and finite limits
exist; hence these environments remain admissible. Finite rational
reward alphabets have a common denominator and satisfy the same
lattice premise. An unrestricted rational reward alphabet does not.

The proposition explains why an argument based on ordering one
agent across several environments cannot establish an impossibility
under the weaker printed definition. It establishes an obstruction
for a different contract. It does not classify every case under
that stronger contract.

\section{Interpretation and further questions}

The positive maps retain a single first-action probability and may
collapse almost all other behavior. That is enough because the target
environment map need preserve only one environment's ordering at a
time. The negative argument identifies a feature that this collapse
cannot erase: a nontrivial stochastic source has intermediate policy
values, whereas a scripted deterministic target environment can make
all deterministic values binary up to a common offset.

Consequently the classification is a result about the expressive
strength of a ranking-preservation definition. It supplies no general
method for compiling a trained policy into an equally capable policy
in another implementation. Requiring computable maps, injectivity,
preservation of actual rewards, a reward normalization, or
cross-environment comparisons would each define a further problem.
These changes should be stated separately before attempting another
classification.

\paragraph{Preparation.}
This is an AI-assisted hand-proof manuscript. The complete four-framework
argument is included, using the versioned source definition. No formal
verification or established historical priority is claimed.

\begin{thebibliography}{9}
\bibitem{AP} S.~A.~Alexander and A.~P.~Pedersen,
\emph{Representation and Invariance in Reinforcement Learning},
\href{https://arxiv.org/abs/2112.07752v5}{arXiv:2112.07752v5},
11 August 2026. Especially Definitions~2 and~6--9 and the concrete
classification section.
\bibitem{Counterexample} Logic Lineages project,
\emph{A counterexample to the transformation classification as written},
working note, 2026.
\href{https://github.com/avg-netizen/alexander-open-problems-results/tree/main/05-rl-transformation-counterexample}{Public note and proof}.
The present manuscript includes all arguments it needs from that note.
\end{thebibliography}
\end{document}
